% Page budget: 0.5 page (W4, W11). Grouped by topic; each paper is described only at the level of its title and venue (or what Sections 2-4 already say).
\section{Related Work}
\label{sec:related}

\emph{Data-center demand response.} Liu et al.\ study how a data center can avoid the coincident peak by shifting workload and using local generation~\cite{liu13}, and Zhou et al.\ compare the primary-response potential of GPU and CPU data centers~\cite{zhou24}. A field test cut the power of a production AI cluster for hours by controlling its workload, without the cloud~\cite{colangelo26}. Closest to our work, Acun et al.\ study the capping and scheduling flexibility of AI data centers for demand response~\cite{acun26dr}; we price cloud offloading against delay and capping on two production traces.

\emph{Carbon-aware computing} uses the same flexibility to lower emissions: in data centers~\cite{radovanovic23}, when scheduling DNN inference on heterogeneous GPUs~\cite{nabavinejad25}, through workload elasticity~\cite{hanafy24}, and when designing data centers~\cite{acun23}. We price a grid request in dollars and treat carbon only as a side check.

\emph{Hybrid-cloud bursting and cost-aware scheduling.} Prior work studies the cost-benefit of extending clusters with cloud resources~\cite{assuncao09}, a cost model for bursting iterative MapReduce jobs~\cite{clemente17}, cost-aware schedulers for hybrid and public clouds~\cite{luo24,chang25}, a broker that runs jobs across clouds~\cite{yang23}, and spot instances under deadlines~\cite{wu24}. These works do not treat bursting as a priced response to a grid request, in which only the hours of the event have value.

\emph{GPU power and energy.} Zeus~\cite{you23} and Perseus~\cite{chung24} study the energy and speed of training on GPUs, Zhao et al.\ measure power capping at HPC scale~\cite{zhao23}, and Patel et al.\ characterise power-management opportunities for large language models in the cloud~\cite{patel24}; DynamoLLM~\cite{stojkovic25dynamo} and TAPAS~\cite{stojkovic25tapas} manage the power of LLM inference clusters. Our cap is a cluster-wide level held for the length of an event.

\emph{Runtime and failure prediction.} Predictions generated from a user's recent jobs can replace user runtime estimates in backfilling~\cite{tsafrir07}. Workload studies of large GPU clusters~\cite{jeon19,hu24,li26} describe the jobs such predictors face, learned predictors for GPU clusters use features of the job and of the user's history~\cite{hu21,hu23}, and learned power predictions guide energy-aware container scheduling~\cite{raith24}. In our runs, how a prediction is read for long-running jobs changed the cost far more than the accuracy differences between the learned predictors (Section~\ref{sec:eval}).
